
\documentclass[12pt]{article}
\usepackage[margin=1in]{geometry}
\usepackage{amsmath,amssymb}
\usepackage{hyperref}
\begin{document}

\section*{Introduction}

The Scalar--Angular--Torsion (SAT) framework arises from a minimal ontological commitment: that the worldlines traced by particles in spacetime diagrams---long treated as purely representational---should instead be considered physical, one-dimensional structures embedded in a four-dimensional manifold. These structures, which we term \emph{filaments}, interact with a propagating three-dimensional surface $\Sigma_t$, interpreted as a resolving time wavefront. The statistical, topological, and geometrical properties of these filament-wavefront interactions constitute the foundational machinery of SAT.

The development of SAT began as a geometric translation of familiar visual tools in physics, particularly those used in relativistic particle kinematics and Feynman diagrammatics. The key postulate is that observed physical properties---mass, charge, spin, and interaction profiles---may emerge from the local and global structure of these filaments and their intersections with $\Sigma_t$. Remarkably, this leads to an internally consistent and increasingly predictive framework that reproduces many structural features of the Standard Model and General Relativity without postulating fields, gauge symmetries, or curvature as primitive.

Crucially, SAT is not an attempt to replace existing frameworks. The Standard Model (SM), General Relativity (GR), and Quantum Mechanics (QM) are treated as boundary constraints---empirically validated, structurally complete within their domains, and serving as essential targets for consistency. The SAT framework does not re-derive these theories from scratch, but instead seeks a geometric substrate from which their formal structures can emerge naturally. Accordingly, SAT takes as foundational the empirical data underpinning SM, GR, and QM---notably particle spectra, coupling constants, and cosmological observables---and constrains itself to reproduce these within a theory defined by two dimensional constants: a filament tension $T$ and a fundamental length scale $\ell$.

What distinguishes SAT is the emergence of coherent structural relationships that extend beyond initial assumptions. For example, the classification of gauge groups as topological binding classes (e.g., $n=2$ for $SU(2)$, $n=3$ for $SU(3)$), the derivation of gauge couplings from filament ensemble densities, and the recovery of mass hierarchies via topological suppression laws ($1/Q$ scaling) all arise from a shared formal basis. These results were not injected, but surfaced organically from a system constructed to model only geometric intersections and their strain responses. The same machinery yields a parameter-free effective field theory for $\Sigma_t$, reproducing the scalar spectral index $n_s \approx 0.965$ and a tensor-scalar ratio $r \approx 0.11$ in alignment with current cosmological data.

SAT remains in an early stage of development, and many derivations are still heuristic or assisted by symbolic computation. Nevertheless, the theory demonstrates a surprising degree of internal consistency, dimensional constraint, and empirical alignment. It offers a conceptual scaffold---not a replacement for known physics, but a possible geometric foundation that connects disparate regimes within a unifying topological and statistical grammar.

The purpose of this manuscript is to formalize the SAT program, summarize its core axioms, derive its principal phenomenological outputs, and assess its limitations and open problems. Where possible, analytic proofs are supplemented by numerical simulations and machine-verified derivations. While SAT may ultimately serve as an intermediate theory---a scaffold rather than a replacement---it provides a coherent new lens through which to examine the structural unity of fundamental physics.

\end{document}
